\documentclass[sigconf,anonymous=false]{acmart}

% --- Draft mode: suppress ACM rights-management block (no DOI/ISBN yet) ---
\settopmatter{printacmref=false}
\renewcommand\footnotetextcopyrightpermission[1]{}
\acmConference[Draft for review]{ACM SIGCITE 2026}{November 12--14, 2026}{Wilmington, NC, USA}
\acmYear{2026}
\copyrightyear{2026}
\acmDOI{}
\acmISBN{}
\acmPrice{}

\usepackage{booktabs}
\usepackage{amsmath}
\usepackage{colortbl}
\usepackage{mdframed}
\usepackage{pifont}

% --- Tighten float/caption spacing to reclaim vertical space (content-neutral) ---
\setlength{\textfloatsep}{3pt plus 1pt minus 2pt}
\setlength{\floatsep}{2pt plus 1pt minus 2pt}
\setlength{\intextsep}{3pt plus 1pt minus 2pt}
\setlength{\abovecaptionskip}{2pt}
\setlength{\belowcaptionskip}{2pt}
\setlength{\abovedisplayskip}{4pt plus 1pt minus 2pt}
\setlength{\belowdisplayskip}{4pt plus 1pt minus 2pt}
\setlength{\abovedisplayshortskip}{2pt plus 1pt minus 1pt}
\setlength{\belowdisplayshortskip}{2pt plus 1pt minus 1pt}

\newcommand{\cmark}{\ding{51}}
\newcommand{\xmark}{\ding{55}}

\definecolor{diagAccent}{HTML}{1F4E5F}
\definecolor{diagAccentSoft}{HTML}{E7EEF1}
\definecolor{flagRed}{HTML}{7A1F2B}
\definecolor{flagRedSoft}{HTML}{F6E9EA}

\mdfdefinestyle{keyinsight}{
  linecolor=diagAccent, linewidth=0.9pt,
  backgroundcolor=diagAccentSoft,
  innerleftmargin=8pt, innerrightmargin=8pt,
  innertopmargin=4pt, innerbottommargin=4pt,
  skipabove=4pt, skipbelow=4pt,
  roundcorner=3pt,
}

\begin{document}

\title{Two Failure Modes in Categorical Target-Leakage Screening: Sparsity Bias and Complete-Case Deletion}

\author{Jaime Cano Mora\~{n}o}
\affiliation{%
  \institution{Illinois Institute of Technology}
  \country{USA}
}
\email{jcanomorano@hawk.iit.edu}

\renewcommand{\shortauthors}{Cano Mora\~{n}o}

\begin{abstract}
Categorical target-leakage screening typically thresholds Cram\'er's V.
We show two independent failure modes on the same statistic, each traceable
to a standard tooling default rather than to $V$ itself.
\emph{Sparsity bias}: the Titanic \texttt{name} column scores $V=0.999$
against survival despite carrying no information, because 1{,}305 of
1{,}307 names are singleton categories; a bias
correction~\cite{bergsma2013bias} collapses this to $\widetilde{V}=0.000$,
while two well-behaved low-cardinality columns barely move
($\texttt{sex}$: $0.529{\to}0.528$; $\texttt{embarked}$: $0.184{\to}0.180$).
\emph{Complete-case deletion}: \texttt{pandas.crosstab}, and the PPS-based
association used by ydata-profiling and Deepchecks'
\texttt{Feature\-Label\-Correlation}, all drop rows with missing feature values
before testing association---the default nearly everywhere. On Titanic's
\texttt{boat} column (a post-hoc leak, 62.9\% missing, informative
missingness), this drops the measured association from $V=0.948$ (full
data) to $V=0.420$ (non-missing subset only), hiding the leak. We fix this
in our own pipeline---encoding missingness as an explicit category, matching
how mutual information already treats it---recovering $V=0.951$ and raising
our unified score from $\mathcal{L}=0.739$ (warning) to $\mathcal{L}=0.925$
(error). A negative replication on NSL-KDD (125{,}973 rows; three
categorical features spanning cardinality 3--70, no missing values)
confirms the sparsity mechanism requires genuine per-category sparsity, not
cardinality alone: raw and corrected $V$ are numerically indistinguishable
there. Neither failure is fixed by tuning a threshold.
\end{abstract}

\ccsdesc[500]{Computing methodologies~Machine learning}
\ccsdesc[300]{Mathematics of computing~Statistical paradigms}
\ccsdesc[300]{Security and privacy~Software security engineering}

\keywords{target leakage, categorical data, Cram\'er's V, bias correction, missing data, data quality}

\maketitle

\begin{mdframed}[style=keyinsight]
\textbf{Key insight.} Two standard defaults---chi-squared bias on sparse
categories, and complete-case deletion of missing feature values---each
break categorical leakage screening in a different, diagnosable way. One is
fixed by a bias-corrected estimator; the other requires treating
missingness itself as a feature. Neither fix addresses the other's failure.
\end{mdframed}

\section{Motivation}

Standard categorical leakage screening thresholds an association statistic,
usually Cram\'er's V, between each feature and the target. Throughout this
paper, \emph{leakage} means \emph{target
leakage}~\cite{kaufman2012leakage}: information about the label that
enters training but would not be available at prediction time---not a
security breach or data exfiltration. We show that on a single canonical
dataset, this statistic fails in two directions for two distinct reasons,
and that both reasons are properties of standard tooling defaults---not of
$V$, and not idiosyncratic to our own implementation. The risk is not
merely academic: near-unique identifier columns and administrative codes
recorded post-outcome are common in exactly the security- and
safety-relevant settings---fraud detection, credit scoring, clinical
triage---where a missed leak is costly.

\section{The Two Failures}
\label{sec:failures}

Table~\ref{tab:main} reports, for four Titanic columns (five rows, since
\texttt{boat} appears both before and after the fix described below),
cardinality, singleton-category count, raw and bias-corrected Cram\'er's
V, normalised mutual information $\widetilde{I}$, performance inflation
$\pi$, the unified score $\mathcal{L}$ computed from each version of the
correlation term ($\mathcal{L}(V) = 0.35V + 0.35\widetilde{I} + 0.30\pi$,
and likewise for $\mathcal{L}(\widetilde{V})$), and the ground-truth
leakage label. Reporting $\mathcal{L}$ under both statistics matters:
under raw $V$, \texttt{name} and \texttt{sex} are nearly indistinguishable
despite one being leaky and the other clean; under corrected
$\widetilde{V}$ they cleanly separate. The two \texttt{boat} rows carry
the same logic in reverse, discussed in Section~\ref{sec:multisignal}.

\begin{table*}[t]
\centering
\caption{Five rows across four Titanic columns. \texttt{name} is a false
  positive under raw $V$; \texttt{boat} is shown both before and after our
  missingness fix (Section~\ref{sec:failures}), so the failure and the
  repair both appear here rather than in prose alone. \texttt{boat}'s
  cardinality differs between rows: pre-fix drops the 823 missing rows
  before counting categories (27); post-fix encodes missingness as its own
  category (28). $^\dagger$Clipped at the estimator's floor, not an exact
  zero.}
\label{tab:main}
\small
\begin{tabular}{lrrrrrrrrl}
\toprule
\rowcolor{diagAccent}
\textcolor{white}{\textbf{Feature}} & \textcolor{white}{Card.} & \textcolor{white}{Sing.} & \textcolor{white}{$V$} & \textcolor{white}{$\widetilde{V}$} & \textcolor{white}{$\widetilde{I}$} & \textcolor{white}{$\pi$} & \textcolor{white}{$\mathcal{L}(V)$} & \textcolor{white}{$\mathcal{L}(\widetilde{V})$} & \textcolor{white}{\textbf{Truth}} \\
\rowcolor{flagRedSoft}
\texttt{name}     & 1307 & 1305 & \textbf{0.999} & 0.000$^\dagger$ & 0.164 & 0.000 & 0.407 & \textbf{0.058} & clean (FP under $V$) \\
\rowcolor{flagRedSoft}
\texttt{boat} (pre-fix)  & 27 & 4 & 0.420 & 0.351 & 1.000 & 0.808 & 0.739 & 0.715 & leaky (FN under $V$) \\
\texttt{boat} (post-fix) & 28 & 4 & 0.951 & 0.940 & 1.000 & 0.808 & \textbf{0.925} & 0.922 & leaky (recovered) \\
\texttt{sex}      & 2    & 0    & 0.529 & 0.528 & 0.265 & 0.424 & 0.405 & 0.405 & clean \\
\texttt{embarked} & 3    & 0    & 0.190 & 0.184 & 0.054 & 0.000 & 0.085 & 0.083 & clean \\
\bottomrule
\end{tabular}
\end{table*}

\paragraph{Why \texttt{name} inflates.}
With 1{,}309 passengers and 1{,}307 distinct names, 1{,}305 names are
singleton categories, and chi-squared-based association statistics are
upward-biased whenever category counts are this sparse relative to sample
size. The bias correction of Bergsma~\cite{bergsma2013bias},
\[
\widetilde{V} = \sqrt{\frac{\max(0,\ \varphi^2 - \frac{(r-1)(c-1)}{n-1})}{\min(\tilde r - 1,\ \tilde c - 1)}},
\quad \tilde r = r - \tfrac{(r-1)^2}{n-1},\ \tilde c = c - \tfrac{(c-1)^2}{n-1},
\]
collapses \texttt{name} to $\widetilde{V}=0.000$---the estimator's clipped
floor (the $\max(0,\cdot)$ term), not a coincidentally exact zero---while
\texttt{sex} and \texttt{embarked} (Table~\ref{tab:main}) barely move,
confirming the correction targets sparsity specifically.

\paragraph{Why \texttt{boat} was missed, and by whom else.}
Lifeboat assignment was recorded only for passengers whose fate was already
known, making \texttt{boat} a textbook post-hoc leak, yet on our original
pipeline it scored only $V=0.420$. Correlation is not blind to this leak
in principle: 62.9\% of \texttt{boat} is missing, and that missingness is
itself almost deterministic of the outcome (Table~\ref{tab:missing})---98.1\%
of passengers with a boat number survived, versus 2.8\% of those without
one. \texttt{pandas.crosstab} excludes rows with a missing value by
default: \texttt{ydata-profiling}'s own Cram\'er's V implementation calls
\texttt{pandas.crosstab} directly, and empirically, Deepchecks' PPS-based
\texttt{FeatureLabelCorrelation} scores \texttt{boat} at essentially zero
(PPS $=0.0006$) with real missing values present, versus $0.946$ once
missingness is an explicit indicator---the same failure, independently
found. This is a shared default across the tooling ecosystem, not an
isolated implementation bug.

\begin{table}[h]
\centering
\caption{Where the original \texttt{boat} signal was lost. This table
  collapses \texttt{boat} to a binary missing/present indicator, isolating
  missingness's own association; $V=0.951$ elsewhere (Table~\ref{tab:main})
  instead keeps all 27 non-missing categories distinct plus a 28th missing
  category, so the two are not directly comparable---both confirm the same
  signal survives, measured two ways.}
\label{tab:missing}
\small
\begin{tabular}{lrr}
\toprule
\rowcolor{diagAccent}
\textcolor{white}{\textbf{Comparison}} & \textcolor{white}{$n$} & \textcolor{white}{$V$} \\
\texttt{boat}.notna() vs.\ survived (full data) & 1309 & \textbf{0.948} \\
\rowcolor{flagRedSoft}
\texttt{boat} vs.\ survived (non-missing only) & 486 & 0.420 \\
\bottomrule
\end{tabular}
\end{table}

We fixed this in our own pipeline: missing feature values are now encoded
as an explicit category rather than dropped, matching how mutual
information already treats them (\texttt{pandas.factorize} assigns missing
values their own code rather than \texttt{NaN}, so our MI and performance
terms never dropped this signal---only correlation did). Table~\ref{tab:main}'s
two \texttt{boat} rows show the result: $\mathcal{L}$ crosses from
\emph{warning} into \emph{error} severity.

\section{Why Multiple Signals Help---and Why One Fix Does Not Cover Both Failures}
\label{sec:multisignal}

Before the fix, $\widetilde{I}$ and $\pi$ already saturated for
\texttt{boat} while $\rho$ did not (\texttt{boat} (pre-fix), Table~\ref{tab:main}:
both $\mathcal{L}(V)$ and $\mathcal{L}(\widetilde{V})$ clear $0.7$ on
$\widetilde{I}$ and $\pi$ alone), because mutual information
estimation~\cite{kraskov2004estimating} never dropped the missing rows to
begin with. This is a real, useful property of combining signals, but it
is not the fix for \texttt{boat}: the two rows show the bias correction
changes almost nothing for it, before the fix
($V{=}0.420\to\widetilde{V}{=}0.351$) or after
($V{=}0.951\to\widetilde{V}{=}0.940$)---sparsity correction and
missingness-aware association are two different remedies for two different
failures, and applying one does not substitute for the other. The weights
$(0.35, 0.35, 0.30)$ are a design choice, not a fitted parameter; we do not
present the weighted sum itself as a research contribution.

\section{Negative Replication on NSL-KDD}
\label{sec:nslkdd}

To test whether the sparsity mechanism generalises, we repeated the raw-vs.-corrected
comparison on NSL-KDD~\cite{tavallaee2009nslkdd} (125{,}973 rows), using its
three native categorical features against a binary attack/normal label
(46.5\% attack rate). Table~\ref{tab:nslkdd} reports the result: negative.

\begin{table}[h]
\centering
\caption{NSL-KDD: raw and corrected $V$ are indistinguishable at every
  cardinality tested---no column here is sparse enough to trigger bias.}
\label{tab:nslkdd}
\small
\begin{tabular}{lrrrr}
\toprule
\rowcolor{diagAccent}
\textcolor{white}{\textbf{Feature}} & \textcolor{white}{Card.} & \textcolor{white}{Sing.} & \textcolor{white}{$V$} & \textcolor{white}{$\widetilde{V}$} \\
\texttt{protocol\_type} & 3  & 0 & 0.282 & 0.282 \\
\texttt{service}        & 70 & 1 & 0.860 & 0.860 \\
\texttt{flag}           & 11 & 0 & 0.775 & 0.775 \\
\bottomrule
\end{tabular}
\end{table}

None of the three columns has any missing values, so this replication
isolates the sparsity mechanism specifically. Even \texttt{service}, at 70
categories, has only 1 singleton among 125{,}973 rows---nothing like
\texttt{name}'s 1{,}305 of 1{,}307---so the correction has nothing to
remove. Despite \texttt{service} scoring as high as $V{=}0.860$, none of
the three features is itself a leak by the same criterion that makes
\texttt{boat} one: \texttt{protocol\_type}, \texttt{service}, and
\texttt{flag} are properties of a connection observable at classification
time, not values recorded after the attack was labelled, so their
association with the outcome is legitimately predictive rather than
post-hoc. NSL-KDD therefore offers no false-positive or false-negative
case here; it confirms instead that the sparsity mechanism tracks
observations-per-category, not cardinality by itself.

\section{Practical Implications}

Any feature that is a discrete, deterministic function of the
outcome---assignment codes, post-event labels, administrative
flags---can evade a thresholded association statistic while remaining
fully visible to mutual information and single-feature performance. We
recommend three practices: (1)~bias-corrected Cram\'er's V on
high-cardinality columns; (2)~missingness indicators screened as features
in their own right, since informative missingness survives complete-case
deletion in every tool we tested; (3)~never rely on a single statistic.
Titanic remains one of the most widely used datasets in introductory ML
instruction, and the screening practice it commonly illustrates is exactly
the one shown insufficient here. The bias-corrected estimator and the
full open-source pipeline (349 automated tests) are available at
\url{https://github.com/jaimecano12/ml-framework}.

\bibliographystyle{ACM-Reference-Format}
\begingroup
\small
\begin{thebibliography}{9}

\bibitem{kaufman2012leakage}
S. Kaufman, S. Rosset, C. Perlich, and O. Stitelman.
2012.
Leakage in data mining: Formulation, detection, and avoidance.
\textit{ACM Transactions on Knowledge Discovery from Data} 6, 4 (2012), 1--21.

\bibitem{bergsma2013bias}
W. Bergsma.
2013.
A bias-correction for Cram\'er's V and Tschuprow's T.
\textit{Journal of the Korean Statistical Society} 42, 3 (2013), 323--328.
\url{https://doi.org/10.1016/j.jkss.2012.10.002}

\bibitem{kraskov2004estimating}
A. Kraskov, H. St\"{o}gbauer, and P. Grassberger.
2004.
Estimating mutual information.
\textit{Physical Review E} 69, 6 (2004), 066138.

\bibitem{tavallaee2009nslkdd}
M. Tavallaee, E. Bagheri, W. Lu, and A. Ghorbani.
2009.
A detailed analysis of the KDD CUP 99 data set.
In \textit{Proc.\ IEEE Symposium on Computational Intelligence for Security and Defense Applications (CISDA)}.

\end{thebibliography}
\endgroup

\end{document}
